\documentclass[11pt,a4paper]{article}
\usepackage[T1]{fontenc}
\usepackage[utf8]{inputenc}
\usepackage{amsmath,amssymb,amsthm}
\usepackage[margin=1in]{geometry}
\usepackage[colorlinks=true,linkcolor=blue,urlcolor=blue]{hyperref}
\theoremstyle{plain}
\newtheorem{theorem}{Theorem}
\newtheorem{lemma}[theorem]{Lemma}
\newtheorem{proposition}[theorem]{Proposition}
\newtheorem{corollary}[theorem]{Corollary}
\theoremstyle{definition}
\newtheorem{remark}[theorem]{Remark}

\title{The empirical recurrence for OEIS A267872, proved from an exact model}
\author{Adrian Perez Fontelles\\ \small Independent researcher}
\date{2 September 2026}
\begin{document}
\maketitle

\begin{abstract}
OEIS A267872 carries an empirical recurrence of order $2$ contributed by Colin Barker. It is true, and it is
decidable rather than empirical. The entry's sequence is given exactly by a certified growth pattern of the automaton's row, from
which $a$ provably satisfies a MONIC linear recurrence of order $S=7$, derived below and not
assumed. The residual of the conjectured recurrence satisfies the same one, so whether it
annihilates $a$ is settled by a finite exact computation. No transfer matrix is involved and
the count is not a walk.
\end{abstract}

\noindent\small 2020 Mathematics Subject Classification. 68Q80, 11B85, 05A15.\normalsize

\section{The conjecture}

OEIS A267872 is ``Number of ON (black) cells in the n-th iteration of the "Rule 237" elementary cellular automaton starting with a single ON (black) cell.''. It has offset $0$ and begins
\[
1, 1, 5, 7, 9, 11, 13, 15,\ \dots
\]
The entry states, as a conjecture contributed by its author and never marked settled:
\begin{quote}
Conjectures from \_Colin Barker\_, Jan 22 2016 and Apr 20 2019: (Start) a(n) = 2*a(n-1)-a(n-2) for n>3.
\end{quote}
As of the ``Last modified'' line on the live entry (Feb 16 2025 08:33:30, revision 23)
this is still recorded as empirical, and nothing on the entry records it as proved.

\section{What is being counted}

The entry counts the ON (black) cells at stage $n$ of the automaton. Nothing is being enumerated over a window, so there is
no digraph and no transfer matrix. What settles the sequence is the shape of the row itself:
the row grows only at its two ends. Writing
$r$ for the residue of $n$ modulo the period $p$, the automaton was run and
\[
w(n+p)\;=\;L_r + w(n) + R_r
\]
was verified for every available $n$ past a
pre-period $n_0$, with $L_r$ and $R_r$ fixed words depending only on $r$. Counting ON cells on
both sides gives
\[
\mathrm{on}(n+p)\;=\;\mathrm{on}(n)+\mathrm{ones}(L_r)+\mathrm{ones}(R_r),
\]
a first-order difference along each class.

The certificate is a theorem and not a fit, for the same reason: the rule is a local map of
radius $1$, so verifying the identity at two consecutive $n$ in a class fixes every three-cell
window at both boundaries, and induction carries it to every later $n$.

For this entry the automaton is rule 237, the period is $p=1$ and the pre-period is $n_0=2$.
The certificate is
\begin{gather*}
r=0:\quad w(n+p) = \varepsilon \cdot w(n) \cdot \texttt{11}
\end{gather*}

\section{The bound}

So
$\mathrm{on}$ is annihilated by $(z^p-1)(z-1)$, and $\mathrm{off}(n)=2n+1-\mathrm{on}(n)$ by
the same polynomial, since the width is linear in $n$; a running total contributes one further
factor $(z-1)$. Allowing the pre-period to raise the numerator's degree gives the monic
annihilator of order $S=7$ used below.

\section{The proof}

\begin{lemma}\label{lem:crit}
Let $A(z)=z^S-\sum_{j<S}\alpha_jz^{j}$ be the monic annihilator derived above and
$q(z)=z^{2}-\sum_ic_iz^{2-i}$ the characteristic polynomial of the conjectured
recurrence. The residual $r(n)=a(n)-\sum_ic_ia(n-i)$ is a linear combination of shifts of $a$,
so $A$ annihilates $r$ as well. Since $A(0)\ne0$ the recurrence it gives runs in both
directions, and $S$ consecutive zeros of $r$ force $r$ to vanish at every later index.
\end{lemma}

\begin{proof}
$A$ annihilates $a$, hence every shift of $a$, hence any linear combination of shifts; that is
$r$. A monic recurrence with nonzero constant term determines each term from the $S$ before it
and each from the $S$ after it, so a run of $S$ zeros propagates.
\end{proof}

\begin{theorem}
\[
a(n)\;=\;2\,a(n-1) - a(n-2)
\]
for every $n>3$.
\end{theorem}

\begin{proof}
The terms $a(0),\dots$ were computed exactly in integer arithmetic from the model, extended by
$A$ where the model itself was not evaluated, and $r(n)$ formed for each index. Those integers
vanish from $n=3+1$ onwards, and the run exceeds $S+2$, which by
Lemma~\ref{lem:crit} settles every larger index at once. The bound is exact: at $n=3$ the recurrence fails on the entry's own published terms, so no larger range holds.
\end{proof}

\section{Verification}

Three checks, each independent of the algebra above.

First, the model reproduces all $63$ terms the entry publishes, exactly, in integer
arithmetic. This is what ties the model to the entry: it is built by reading the entry's
English, and a misreading gives different counts.

Second, the conjectured recurrence was evaluated directly on those published terms, with no
model involved, and holds wherever the proved range and the data overlap.

Third, the derived annihilator $A$ was itself tested on terms it had not been given: the model
was evaluated beyond the $S$ values that determine it and $A$ reproduced them. A bound that is
too small fails this test, which is why it is run.

\begin{thebibliography}{9}
\bibitem{oeis} The OEIS Foundation, \emph{The On-Line Encyclopedia of Integer
Sequences}, \url{https://oeis.org}, 2026. Sequence A267872.
\bibitem{beckrobins} M.~Beck and S.~Robins, \emph{Computing the Continuous Discretely},
2nd ed., Springer, 2015. (Ehrhart quasi-polynomials and their periods.)
\bibitem{stanley} R.~P.~Stanley, \emph{Enumerative Combinatorics}, Volume 1, 2nd ed.,
Cambridge University Press, 2011.
\bibitem{ds} F.~Diamond and J.~Shurman, \emph{A First Course in Modular Forms},
Springer GTM 228, 2005.
\end{thebibliography}

\end{document}
